\subsection{枝の長さ2}

\begin{definition}
  \label{def:adjExcl}
  \lean{adjExcl}
  \leanok
  \uses{}
  $a, b, u \in V$とする．
  $u$を除いたグラフにおいて，$a$と$b$が隣接しているとき，$a \sim_u b$と書く：
  \begin{equation}
    a \sim_u b
    :\iff a \ne b,\ a \ne u,\ b \ne u,\ C_{ab} \ne 0.
  \end{equation}
\end{definition}

\begin{definition}
  \label{def:reachExcl}
  \lean{reachExcl}
  \leanok
  \uses{}
  $u, v, w \in V$とする．
  $u$を通らずに$v$から$w$まで到達可能なとき，$v \leadsto_u w$と書く．
\end{definition}

\begin{definition}
  \label{def:branchSet}
  \lean{branchSet}
  \leanok
  $u, v \in V$とする．
  集合$B_u(v) \subseteq V$を以下で定める：
  \begin{equation}
    B_u(v) := \set{w}{v \leadsto_u w}.
  \end{equation}
  つまり，頂点$u$を削除した後に$v$と繋がっている頂点たちの集合である．
\end{definition}

\begin{lemma}
  \label{lem:branchFinset_refl}
  \lean{branchFinset_refl}
  \leanok
  \uses{}
  $u, v \in V$とする．
  このとき，$v \in B_u(v)$が成り立つ．
\end{lemma}

\begin{lemma}
  \label{lem:branch_connected}
  \lean{branch_connected}
  \leanok
  \uses{}
  $u, v \in V$とする．
  このとき，任意の$w \in B_u(v)$に対し，$B_u(v)$内のみを通って$v$から$w$まで到達可能である：
  \begin{equation}
    v \reach{B_u(v)} w.
  \end{equation}
\end{lemma}

\begin{proof}
  頂点に関する帰納法で示す．

  $w = v$のときはよい．

  $w = a$までは正しい，すなわち$v \reach{B_u(v)} a$とする．
  $v \leadsto_u a,\ a \sim_u b$とし，$v \reach{B_u(v)} b$を示す．
  $v \leadsto_u a$より$a \in B_u(v)$である．
  さらに，$a \sim_u b$であるから$b \in B_u(v)$である．
  したがって，$a \reach{B_u(v)} b$であるから$v \reach{B_u(v)} b$である．
\end{proof}

\begin{lemma}
  \label{lem:branch_marking}
  \lean{branch_marking}
  \leanok
  \uses{lem:branch_energy_strong, lem:branch_connected}
  $C$はGCMとし，$u, v \in V$とする．
  このとき，$y := \transpose{(y_1, \ldots, y_n)} \in \R^n$で以下を満たすものが存在する：
  \begin{enumerate}[label = (y\arabic*)]
    \item $\forall i \in V,\ 0 \le y_i$.
    \item $\forall i \in V,\ i \notin B_u(v) \implies y_i = 0$.
    \setcounter{enumi}{3}
    \item $\displaystyle \frac{\# B_u(v)}{\# B_u(v) + 1} \le 2 y_v - 2 \sum_{i \in B_u(v)} y_i ^ 2 + \sum_{i \in B_u(v)} \sum_{j \in B_u(v)} \bm{1}_{i \sim j}\ y_i y_j$.
  \end{enumerate}
\end{lemma}

\begin{proof}
  $v \in B_u(v)$である．
  補題~\ref{lem:branch_connected}より，補題~\ref{lem:branch_energy_strong}の前提条件は，$K := B_u(v),\ r := v$とすれば満たされる．
  このとき，補題~\ref{lem:branch_energy_strong}の主張は，まさに示したい(1), (2), (4)である．
\end{proof}

\begin{lemma}
  \label{lem:reachExcl_ne_u}
  \lean{reachExcl_ne_u}
  \leanok
  \uses{}
  $u, v, w \in V$とし，$v \leadsto_u w,\ v \ne u$とする．
  このとき，$w \ne u$である．
\end{lemma}

\begin{proof}
  $v$から$w$への道の長さに関する帰納法による．

  $w = v$のときはよい．
  $w = a$までは正しい，すなわち$a \ne u$とする．
  $v \leadsto_u a,\ a \sim_u b$とし，$b \ne u$を示す．
  $a \sim_u b$より従う．
\end{proof}

\begin{lemma}
  \label{lem:u_notMem_branch}
  \lean{u_notMem_branch}
  \leanok
  \uses{lem:reachExcl_ne_u}
  $u, v \in V$とし，$v \ne u$とする．
  このとき，$u \notin B_u(v)$である．
\end{lemma}

\begin{proof}
  $v \not \leadsto_u u$を示せばよい．
  $v \leadsto_u u$と仮定すると，補題~\ref{lem:reachExcl_ne_u}より$u \ne u$となり，矛盾する．
\end{proof}

\begin{lemma}
  \label{lem:neighborSet_eq_triple}
  \lean{neighborSet_eq_triple}
  \leanok
  \uses{}
  $u \in V,\ \deg(u) = 3$とし，$v_1, v_2, v_3 \in N(u)$は相異なるとする．
  このとき，以下が成り立つ：
  \begin{equation}
    N(u) = \{v_1, v_2, v_3\}.
  \end{equation}
\end{lemma}

\begin{proof}
  $\{v_1, v_2, v_3\} \subseteq N(u)$であるから，$\# N(u) \le \# \{v_1, v_2, v_3\}$を示せばよい．
  $\# \{v_1, v_2, v_3\} = 3$である．
  また，$\# N(u) = \deg(u) = 3$であるから成り立つ．
\end{proof}

\begin{lemma}
  \label{lem:form_le_Gadj}
  \lean{form_le_Gadj}
  \leanok
  \uses{}
  $C$はGCMとする．
  $x := \transpose{(x_1, \ldots, x_n)} \in \R^n$とし，任意の$i$に対し，$0 \le x_i$であるとする．
  このとき，以下が成り立つ：
  \begin{equation}
    \transpose{x} S(C) x
    \le 2 \sum_{i \in I} x_i ^ 2 - \sum_{i \in I} \sum_{j \in I} \bm{1}_{i \sim j}\ x_i x_j.
  \end{equation}
\end{lemma}

\begin{proof}
  二次形式を展開する：
  \begin{equation}
    \transpose{x} S(C) x
    = \sum_{i \in I} x_i \left( \sum_{j \in I} S(C)_{ij} x_j \right).
  \end{equation}
  また，任意の$i, j$に対し，次が成り立つことは容易に確かめられる($i = j,\ i \sim j$かどうかの4パターンをそれぞれ調べる)：
  \begin{equation}
    x_i S(C)_{ij} x_j
    \le \bm{1}_{i = j}\ 2 x_i x_j - \bm{1}_{i \sim j}\ x_i x_j.
  \end{equation}
  よって，各項ごとに成り立つので，和を取れば主張が従う．
\end{proof}

\begin{lemma}
  \label{lem:only_neighbor_in_branch}
  \lean{only_neighbor_in_branch}
  \leanok
  \uses{}
  $u, v_1, v_2, v_3 \in V$とし，$N(u) = \{v_1, v_2, v_3\}$とする．
  $B_u(v_1) \cap B_u(v_2) = B_u(v_1) \cap B_u(v_3) = \emptyset$とする．
  また，$i \in V$とし，$i \in B_u(v_1),\ u \sim i$とする．
  このとき，$i = v_1$である．
\end{lemma}

\begin{proof}
  $u \sim i$より$i \in N(u)$であるから，$i$は$v_1, v_2, v_3$のいずれかである．
  また，$v_2 \in B_u(v_2),\ v_3 \in B_u(v_3)$である．
  これらは$B_u(v_1)$と交わらないので，$i = v_1$がわかる．
\end{proof}

\begin{lemma}
  \label{lem:branch_no_cross_edge}
  \lean{branch_no_cross_edge}
  \leanok
  \uses{lem:reachExcl_ne_u}
  $u, v, v' \in V$とし，$v \ne u,\ v' \ne u,\ B_u(v) \cap B_u(v') = \emptyset$であるとする．
  また，$i \in B_u(v),\ j \in B_u(v')$とする．
  このとき，$i \not \sim j$である．
\end{lemma}

\begin{proof}
  背理法で示す．
  $i \sim j$と仮定する．
  補題の仮定より$v \leadsto_u i,\ v' \leadsto_u j$である．
  よって，補題~\ref{lem:reachExcl_ne_u}より$i \ne u,\ j \ne u$である．
  したがって，$v \leadsto_u j$である．
  これは$j \in B_u(v)$を意味するが，$B_u(v) \cap B_u(v') = \emptyset$に矛盾する．
\end{proof}

\begin{lemma}
  \label{lem:assembly_sq}
  \lean{assembly_sq}
  \leanok
  \uses{}
  $u, v_1, v_2, v_3 \in V$とし，$u$は$B_u(v_1), B_u(v_2), B_u(v_3)$にいずれにも属さないとする．
  また，$y_1, y_2, y_3 \in \R^n$とし，任意の$i \in I$に対し，$i \notin B_u(v_j)$ならば$(y_j)_i = 0$であるとする．
  さらに，$B_u(v_1), B_u(v_2), B_u(v_3)$は互いに交わらないとする．
  このとき，以下が成り立つ：
  \begin{equation}
    \sum_{i \in I} g(i) ^ 2
    = 1 + \sum_{i \in B_u(v_1)} (y_1)_i ^ 2 + \sum_{i \in B_u(v_2)} (y_2)_i ^ 2 + \sum_{i \in B_u(v_3)} (y_3)_i ^ 2.
  \end{equation}
  ただし，
  \begin{equation}
    g(i) = \begin{cases}
      1, & i = u, \\
      (y_1)_i + (y_2)_i + (y_3)_i, & \textrm{otherwise}.
    \end{cases}
  \end{equation}
  である．
\end{lemma}

\begin{proof}
  任意の$i$対し，以下が成り立つ：
  \begin{equation}
    [(y_1)_i + (y_2)_i + (y_3)_i] ^ 2 = (y_1)_i ^ 2 + (y_2)_i ^ 2 + (y_3)_i ^ 2.
  \end{equation}
  なぜなら，例えば$i \in B_u(v_1)$のとき，$i \notin B_u(v_2), B_u(v_3)$であるから，$(y_2)_i = (y_3)_i = 0$となるからである．
  よって，
  \begin{align}
    \sum_{i \in I} g(i) ^ 2
    &= 1 + \sum_{i \in I} [(y_1)_i + (y_2)_i + (y_3)_i] ^ 2 - [(y_1)_u + (y_2)_u + (y_3)_u] ^ 2 \\
    &= 1 + \sum_{i \in I} [(y_1)_i ^ 2 + (y_2)_i ^ 2 + (y_3)_i ^ 2] - [(y_1)_u + (y_2)_u + (y_3)_u] ^ 2 \\
  \end{align}
  となる．
  さらに，
  \begin{equation}
    \sum_{i \in I} (y_j)_i ^ 2 = \sum_{i \in B_u(v_j)} (y_j)_i ^ 2,\qquad (y_j)_u = 0
  \end{equation}
  であるから，補題の主張は成り立つ．
\end{proof}

\begin{lemma}
  \label{lem:assembly_cross}
  \lean{assembly_cross}
  \leanok
  \uses{lem:branchFinset_refl}
  $C$はGCMであるとする．
  $u, v_1, v_2, v_3 \in V$とし，$u$は$B_u(v_1), B_u(v_2), B_u(v_3)$にいずれにも属さないとする．
  また，$y_1, y_2, y_3 \in \R^n$とし，任意の$i \in I$に対し，$i \notin B_u(v_j)$ならば$(y_j)_i = 0$であるとする．
  さらに，$B_u(v_1), B_u(v_2), B_u(v_3)$は互いに交わらないとする．
  加えて，$N(u) = \{v_1, v_2, v_3\}$とする．
  このとき，以下が成り立つ：
  \begin{equation}
    \begin{split}
      \sum_{i \in I} \sum_{j \in I} \bm{1}_{i \sim j}\ g(i) g(j)
    {}={}\quad &\sum_{i \in B_u(v_1)} \sum_{j \in B_u(v_1)} \bm{1}_{i \sim j}\ (y_1)_i (y_1)_j \\
    {}+{} &\sum_{i \in B_u(v_2)} \sum_{j \in B_u(v_2)} \bm{1}_{i \sim j}\ (y_2)_i (y_2)_j \\
    {}+{} &\sum_{i \in B_u(v_3)} \sum_{j \in B_u(v_3)} \bm{1}_{i \sim j}\ (y_3)_i (y_3)_j \\
    {}+{} & 2 [(y_1)_{v_1} + (y_2)_{v_2} + (y_3)_{v_3}].
    \end{split}
  \end{equation}
  ただし，
  \begin{equation}
    g(i) = \begin{cases}
      1, & i = u, \\
      (y_1)_i + (y_2)_i + (y_3)_i, & \textrm{otherwise}.
    \end{cases}
  \end{equation}
  である．
\end{lemma}

\begin{proof}
  補題~\ref{lem:assembly_sq}の証明中と同様に，任意の$i, j$に対し，以下が成り立つ：
  \begin{equation}
    \begin{split}
      &\bm{1}_{i \sim j}\ [(y_1)_i + (y_2)_i + (y_3)_i] [(y_1)_j + (y_2)_j + (y_3)_j] \\
    ={} &\bm{1}_{i \sim j}\ (y_1)_i (y_1)_j
    + \bm{1}_{i \sim j}\ (y_2)_i (y_2)_j
    + \bm{1}_{i \sim j}\ (y_3)_i (y_3)_j.
    \end{split}
  \end{equation}
  よって，$g(i)$を展開した後これを適用すると
  \begin{alignat}{2}
    \sum_{i \in I} \sum_{j \in I} \bm{1}_{i \sim j}\ g(i) g(j)
    & = &&\sum_{i \in V \setminus \{u\}} \sum_{j \in V \setminus \{u\}} \bm{1}_{i \sim j}\ [(y_1)_i + (y_2)_i + (y_3)_i] [(y_1)_j + (y_2)_j + (y_3)_j] \\
      &&& +{} \sum_{i \in V \setminus \{u\}} \bm{1}_{u \sim i}\ [(y_1)_i + (y_2)_i + (y_3)_i]
      + \sum_{j \in V \setminus \{u\}} \bm{1}_{j \sim u}\ [(y_1)_j + (y_2)_j + (y_3)_j] \\
    & = &&\sum_{i \in V \setminus \{u\}} \sum_{j \in V \setminus \{u\}} \bm{1}_{i \sim j}\ (y_1)_i (y_1)_j \\
      &&&+{} \sum_{i \in V \setminus \{u\}} \sum_{j \in V \setminus \{u\}} \bm{1}_{i \sim j}\ (y_2)_i (y_2)_j \\
      &&&+{} \sum_{i \in V \setminus \{u\}} \sum_{j \in V \setminus \{u\}} \bm{1}_{i \sim j}\ (y_3)_i (y_3)_j \\
      &&&+{} 2 \sum_{i \in V \setminus \{u\}} \bm{1}_{u \sim i}\ [(y_1)_i + (y_2)_i + (y_3)_i]
  \end{alignat}
  となる．
  したがって，あとはこの式と補題の主張の各項が等しいことを示せばよい．
  \begin{equation}
    \sum_{i \in V \setminus \{u\}} \sum_{j \in V \setminus \{u\}} \bm{1}_{i \sim j}\ (y_1)_i (y_1)_j
    = \sum_{i \in B_u(v_1)} \sum_{j \in B_u(v_1)} \bm{1}_{i \sim j}\ (y_1)_i (y_1)_j
  \end{equation}
  を示す．
  $x \in B_u(v_1)$をとる．
  $x = u$と仮定すると，$x \in B_x(v_1)$であるが，補題の仮定より$x \notin B_x(v_1)$であるから矛盾する．
  よって，$x \ne u$であるから$B_u(v_1) \subseteq V \setminus \{u\}$である．
  したがって，以下の2つを示せばよい：
  \begin{gather}
    \sum_{i \in B_u(v_1)} \sum_{j \in V \setminus \{u\}} \bm{1}_{i \sim j}\ (y_1)_i (y_1)_j
    = \sum_{i \in B_u(v_1)} \sum_{j \in B_u(v_1)} \bm{1}_{i \sim j}\ (y_1)_i (y_1)_j, \\
    \forall x \in V \setminus \{u\},\ x \notin B_u(v_1) \implies \sum_{j \in V \setminus \{u\}} \bm{1}_{i \sim j}\ (y_1)_i (y_1)_j = 0.
  \end{gather}
  前者に再び$B_u(v_1) \subseteq V \setminus \{u\}$を適用すれば，示すべきは以下の3つになる：
  \begin{gather}
    \sum_{i \in B_u(v_1)} \sum_{j \in B_u(v_1)} \bm{1}_{i \sim j}\ (y_1)_i (y_1)_j
    = \sum_{i \in B_u(v_1)} \sum_{j \in B_u(v_1)} \bm{1}_{i \sim j}\ (y_1)_i (y_1)_j, \\
    \forall x \in V \setminus \{u\},\ x \notin B_u(v_1) \implies \sum_{j \in V \setminus \{u\}} \bm{1}_{x \sim j}\ (y_1)_x (y_1)_j = 0, \\
    \forall x \in V \setminus \{u\},\ x \notin B_u(v_1) \implies \bm{1}_{i \sim x}\ (y_1)_i (y_1)_x = 0.
  \end{gather}
  1つ目は明らかである．
  2, 3つ目は，補題の仮定より$x \notin B_u(v_1)$のとき，$(y_1)_x = 0$であるから成り立つ．
  これで，$y_1$に関する項が等しいことが示せた．
  $y_2, y_3$に関する項も全く同様に示せる．

  最後に，
  \begin{equation}
    2 \sum_{i \in V \setminus \{u\}} \bm{1}_{u \sim i}\ [(y_1)_i + (y_2)_i + (y_3)_i]
    = 2 [(y_1)_{v_1} + (y_2)_{v_2} + (y_3)_{v_3}]
  \end{equation}
  を示す．
  $\{v_1, v_2, v_3\} \subseteq V \setminus \{u\}$であるから，以下の2つを示せば良い：
  \begin{gather}
    \sum_{i \in \{v_1, v_2, v_3\}} \bm{1}_{u \sim i}\ [(y_1)_i + (y_2)_i + (y_3)_i]
    = [(y_1)_{v_1} + (y_2)_{v_2} + (y_3)_{v_3}], \\
    \forall x \in V \setminus \{u\},\ x \notin \{v_1, v_2, v_3\} \implies \bm{1}_{u \sim x}\ [(y_1)_x + (y_2)_x + (y_3)_x] = 0.
  \end{gather}
  前者の等式を示す．
  補題~\ref{lem:branchFinset_refl}と$B_u(v_1), B_u(v_2), B_u(v_3)$は互いに交わらないことより，$i \ne j$に関して以下がすべて成り立つ：
  \begin{equation}
    v_i \in B_u(v_i),\qquad v_i \notin B_u(v_j), \qquad v_i \ne v_j.
  \end{equation}
  よって，$(y_j)_{v_i} = 0$であるから，左辺の全9項のうち生き残るのは$(y_i)_{v_i}$の3項だけである．
  後者は，$N(u) = \{v_1, v_2, v_3\}$より，$\bm{1}_{u \sim x} = 0$となるから成り立つ．
\end{proof}

\begin{lemma}
  \label{lem:no_posDef_of_branch_path}
  \lean{no_posDef_of_branch_path}
  \leanok
  \uses{lem:no_cycle}
  $C$はGCMであるとする．
  $u, v, v' \in V$とし，$v \ne v',\ u \sim v,\ u \sim v',\ v \leadsto_u v'$であるとする．
  このとき，$S(C)$は正定値でない．
\end{lemma}

\begin{proof}
  頂点の列$s := (s_1, s_2, \ldots, s_k)$で以下をすべて満たすものが存在することを示す($i \ne j$)：
  \begin{equation}
    s_1 = v,\quad s_k = v',\quad s_i \ne s_j,\quad s_i \sim s_{i+1},\quad s_i \ne u.
  \end{equation}
  $v \leadsto_u v'$であるから，$u$を通らない$v$から$v'$への道$p$が存在する．
  もし$p$に同じ頂点が複数回現れるなら，その間にできる閉じた部分を除去すれば，$p$は単順路にできる．
  $s$は，$p$を構成する頂点を順に並べたものとする．
  すると，$s$が満たすべき上の5条件はすべて従う．

  $L$を$s$の先頭に$u$を加えた列$L = (s_0 := u, s_1, s_2, \ldots, s_k)$とする．
  $u, v, v'$は相異なるから，$3 \le |L|$である．
  また，$s$に重複はなく，$u \notin s$であるから，$L$にも重複がない．

  $L$は巡回的に隣接していること，すなわち任意の$i = 0, 1, \ldots, k$に対し，$j = i + 1 \mod |L|$として$C_{s_i s_j} \ne 0$であることを示す．
  $i = 0, k$のとき，補題の仮定から$C_{uv} \ne 0,\ C_{v'u} \ne 0$が従う．
  $i = 1, \ldots, k-1$のとき，$u$に関係しないので，$s$の性質から従う．

  よって，グラフ内にサイクルがある．
  したがって，定理~\ref{thm:no_cycle}より，$S(C)$は正定値でない．
\end{proof}

\begin{lemma}
  \label{lem:branches_disjoint}
  \lean{branches_disjoint}
  \leanok
  \uses{lem:no_posDef_of_branch_path}
  $C$はGCMで，$S(C)$は正定値であるとする．
  $u, v, v' \in V$とし，$v \ne v'$で$v, v' \in N(u)$とする．
  このとき，$B_u(v) \cap B_u(v') = \emptyset$である．
\end{lemma}

\begin{proof}
  $w \in B_u(v) \cap B_u(v')$として矛盾を導く．
  $v \leadsto_u w$かつ$v' \leadsto_u w$であるから$v \leadsto_u v'$である．
  すると，補題~\ref{lem:no_posDef_of_branch_path}より$S(C)$は正定値でないが，これは矛盾する．
\end{proof}

\begin{lemma}
  \label{lem:schur_ineq}
  \lean{schur_ineq}
  \leanok
  \uses{lem:neighborSet_eq_triple, lem:branch_marking, lem:u_notMem_branch, lem:branches_disjoint, lem:form_le_Gadj, lem:assembly_sq, lem:assembly_cross}
  $C$はGCMで，$S(C)$は正定値であるとする．
  $u \in V$とし，$\deg(u) = 3$とする．
  $v_1, v_2, v_3 \in N(u)$は相異なるとする．
  このとき，次の不等式が成り立つ：
  \begin{equation}
    \frac{\# B_u(v_1)}{\# B_u(v_1) + 1} + \frac{\# B_u(v_2)}{\# B_u(v_2) + 1} + \frac{\# B_u(v_3)}{\# B_u(v_3) + 1}
    < 2.
  \end{equation}
\end{lemma}

\begin{proof}
  $\textrm{(左辺)} \ge 2$として矛盾を導く．

  補題~\ref{lem:neighborSet_eq_triple}より，$N(u) = \{v_1, v_2, v_3\}$である．
  隣接の定義より，$v_1, v_2, v_3 \ne u$である．
  補題~\ref{lem:u_notMem_branch}より，$u \notin B_u(v_1), B_u(v_2), B_u(v_3)$である．
  補題~\ref{lem:branches_disjoint}より，$B_u(v_1), B_u(v_2), B_u(v_3)$は互いに交わらない．

  補題~\ref{lem:branch_marking}より，$y_k \in \R^n$で以下を満たすものがとれる($k = 1, 2, 3$)：
  \begin{enumerate}[label = (y\arabic*)]
    \item $\forall i \in V,\ 0 \le (y_k)_i$.
    \item $\forall i \in V,\ i \notin B_u(v_k) \implies (y_k)_i = 0$.
    \setcounter{enumi}{3}
    \item $\displaystyle \frac{\# B_u(v_k)}{\# B_u(v_k) + 1} \le 2 (y_k)_{v_k} - 2 \sum_{i \in B_u(v_k)} (y_k)_i ^ 2 + \sum_{i \in B_u(v_k)} \sum_{j \in B_u(v_k)} \bm{1}_{i \sim j}\ (y_k)_i (y_k)_j$. \label{enum:y4'}
  \end{enumerate}

  \begin{equation}
    g(i) = \begin{cases}
      1, & i = u, \\
      (y_1)_i + (y_2)_i + (y_3)_i, & \textrm{otherwise}.
    \end{cases}
  \end{equation}
  とする．
  $g(u) \ne 0$であるから，$g \ne 0$である．
  よって，この$g$は$\bm{0}$でない$\R^n$のベクトルだと思えるので，$S(C)$の正定値性より
  \begin{equation}
    0 < \transpose{g} S(C) g
  \end{equation}
  である．
  また，任意の$i \in I$に対して，
  \begin{equation}
    0 \le g(i)
  \end{equation}
  であることは簡単にわかる．
  したがって，補題~\ref{lem:form_le_Gadj}より，
  \begin{equation} \label{eq:hle}
    \transpose{g} S(C) g
    \le 2 \sum_{i \in I} g_i ^ 2 - \sum_{i \in I} \sum_{j \in I} \bm{1}_{i \sim j}\ g_i g_j.
  \end{equation}
  である．

  補題~\ref{lem:assembly_sq}より，
  \begin{equation} \label{eq:hsq}
    \sum_{i \in I} g(i) ^ 2
    = 1 + \sum_{i \in B_u(v_1)} (y_1)_i ^ 2 + \sum_{i \in B_u(v_2)} (y_2)_i ^ 2 + \sum_{i \in B_u(v_3)} (y_3)_i ^ 2.
  \end{equation}
  である．
  さらに，補題~\ref{lem:assembly_cross}より，
  \begin{equation} \label{eq:hcr}
    \begin{split}
      \sum_{i \in I} \sum_{j \in I} \bm{1}_{i \sim j}\ g(i) g(j)
    {}={}\quad &\sum_{i \in B_u(v_1)} \sum_{j \in B_u(v_1)} \bm{1}_{i \sim j}\ (y_1)_i (y_1)_j \\
    {}+{} &\sum_{i \in B_u(v_2)} \sum_{j \in B_u(v_2)} \bm{1}_{i \sim j}\ (y_2)_i (y_2)_j \\
    {}+{} &\sum_{i \in B_u(v_3)} \sum_{j \in B_u(v_3)} \bm{1}_{i \sim j}\ (y_3)_i (y_3)_j \\
    {}+{} & 2 [(y_1)_{v_1} + (y_2)_{v_2} + (y_3)_{v_3}].
    \end{split}
  \end{equation}
  である．

  式~\eqref{eq:hsq}, \eqref{eq:hcr}を式\eqref{eq:hle}に代入し，\ref{enum:y4'}を適用すると，
  \begin{equation}
    \begin{split}
      \transpose{g} S(C) g
    {}\le{} &2 \left( 1 + \sum_{i \in B_u(v_1)} (y_1)_i ^ 2 + \sum_{i \in B_u(v_2)} (y_2)_i ^ 2 + \sum_{i \in B_u(v_3)} (y_3)_i ^ 2 \right) \\
    &{}-{} \sum_{i \in B_u(v_1)} \sum_{j \in B_u(v_1)} \bm{1}_{i \sim j}\ (y_1)_i (y_1)_j \\
    &{}-{} \sum_{i \in B_u(v_2)} \sum_{j \in B_u(v_2)} \bm{1}_{i \sim j}\ (y_2)_i (y_2)_j \\
    &{}-{} \sum_{i \in B_u(v_3)} \sum_{j \in B_u(v_3)} \bm{1}_{i \sim j}\ (y_3)_i (y_3)_j \\
    &{}-{}  2 [(y_1)_{v_1} + (y_2)_{v_2} + (y_3)_{v_3}]. \\
    {}={}& 2 - \sum_{k = 1}^3 \left( 2 (y_k)_{v_k} - 2 \sum_{i \in B_u(v_k)} (y_k)_i ^ 2 + \sum_{i \in B_u(v_k)} \sum_{j \in B_u(v_k)} \bm{1}_{i \sim j}\ (y_k)_i (y_k)_j \right) \\
    {}\le{}& 2 - \sum_{k = 1}^3 \frac{\# B_u(v_k)}{\# B_u(v_k) + 1}
    \end{split}
  \end{equation}
  となる．
  背理法の仮定より，$\dsum_{k = 1}^3 \frac{\# B_u(v_k)}{\# B_u(v_k) + 1} \ge 2$であるから，結局$\transpose{g} S(C) g \le 0$を得る．
  しかし，これは$S(C)$が正定値であることに矛盾する．
\end{proof}

\begin{theorem}
  \label{thm:branchSize_inequality}
  \lean{branchSize_inequality}
  \leanok
  \uses{lem:schur_ineq}
  $C$はGCMで，$S(C)$は正定値であるとする．
  $u \in V$とし，$\deg(u) = 3$とする．
  $v_1, v_2, v_3 \in N(u)$は相異なるとする．
  このとき，次の不等式が成り立つ：
  \begin{equation}
    \frac{1}{\# B_u(v_1) + 1} + \frac{1}{\# B_u(v_2) + 1} + \frac{1}{\# B_u(v_3) + 1}
    > 1.
  \end{equation}
\end{theorem}

\begin{proof}
  補題~\ref{lem:schur_ineq}の結果を整理すれば得られる．
\end{proof}

\begin{lemma}
  \label{lem:reciprocal_sum_gt_one_bounds}
  \lean{reciprocal_sum_gt_one_bounds}
  \leanok
  \uses{}
  $a, b, c \in \N$とし，$0 < a,\ a \le b \le c$であるとする．
  また，次の不等式が成り立つとする：
  \begin{equation}
    \frac{1}{a + 1} + \frac{1}{b + 1} + \frac{1}{c + 1} > 1.
  \end{equation}
  このとき，$a = 1$であり，$b = 1$または$b = 2$かつ$c \le 4$である．
\end{lemma}

\begin{proof}
  丁寧に場合分けすればわかる．
\end{proof}

\begin{theorem}
  \label{thm:branchSize_bounds}
  \lean{branchSize_bounds}
  \leanok
  \uses{thm:branchSize_inequality, lem:branchFinset_refl, lem:reciprocal_sum_gt_one_bounds}
  $C$はGCMで，$S(C)$は正定値であるとする．
  $u \in V$とし，$\deg(u) = 3$とする．
  $v_1, v_2, v_3 \in N(u)$は相異なるとする．
  また，$\# B_u(v_1) \le \# B_u(v_2) \le \# B_u(v_3)$とする．
  このとき，$\# B_u(v_1) = 1$であり，$\# B_u(v_2) = 1$または$\# B_u(v_2) = 2$かつ$\# B_u(v_3) \le 4$である．
\end{theorem}

\begin{proof}
  定理~\ref{thm:branchSize_inequality}の不等式が成り立つ．
  また，補題~\ref{lem:branchFinset_refl}より，$v_1 \in B_u(v_1)$であるから$0 < \# B_u(v_1)$である．
  よって，補題~\ref{lem:reciprocal_sum_gt_one_bounds}が適用でき，主張が従う．
\end{proof}